\documentclass[11pt,a4paper]{article}
\usepackage[utf8]{inputenc}
\usepackage{amsmath,amssymb}
\usepackage{graphicx}
\usepackage{hyperref}
\usepackage{geometry}
\usepackage{booktabs}
\geometry{margin=1in}

\title{Improving Cranfield Information Retrieval via Modular Lexical and Semantic Pipelines}
\author{Project Report}
\date{May 2026}

\begin{document}
\maketitle

\begin{abstract}
This project investigates the limitations of classical lexical retrieval and the degree to which semantic and context-aware methodologies can improve information retrieval on the highly technical Cranfield dataset. We implement an end-to-end modular pipeline partitioned into Query Processing, Document Retrieval, and Ranking blocks, enabling systematic formulation of 72 hybrid configurations. We experiment with advanced lexical techniques (IDF-based query reduction, Unordered Local-Context Bag-of-Words), explicit and latent semantic matching (LSA, ESA), sequence embeddings (WordMover's Distance, SIF, Word2Vec Soft Cosine), and dynamically thresholded semantic query expansion. Our empirical findings demonstrate that while neural dense embeddings dilute aeronautical specificity, Latent Semantic Analysis provides a $+22.4\%$ gain in MAP over the unigram TF-IDF baseline. The optimal configuration merges precise term matching with latent semantic associations, achieving a MAP@10 of 0.3795 and demonstrating that robust query representation combined with SVD-based ranking effectively mitigates vocabulary mismatch while honoring domain-specific precision.
\end{abstract}

% --------------------------------------------------
% 1. Introduction
% --------------------------------------------------
\section{Introduction}
Information retrieval within specialized domains such as aeronautics (represented by the Cranfield collection) presents strict requirements. Standard lexical models like unigram TF-IDF provide strong baseline metrics by matching exact noun phrases and technical specifications. However, they remain highly susceptible to vocabulary mismatch—failing to capture meaning when users queries employ synonyms or paraphrase. 

Rather than adopting a monolithic approach, we structured the project as a multi-stage modular pipeline (\textit{merging} architecture). This pipeline decouples indexing from retrieval and ranking, providing an experimental framework to rigorously test varying hypotheses regarding query reduction, context injection, and latent semantic projection. This report details the theoretical motivations, mathematical formulations, and engineering implementations defining our proposed methodology.

% --------------------------------------------------
% 2. Pipeline Architecture
% --------------------------------------------------
\section{Pipeline Architecture}
To ensure reproducibility and isolate the effect of individual methodologies, we engineered a three-tier pipeline:
\begin{enumerate}
    \item \textbf{Block 1: Query Processing:} Manages tokenization, vectorization, reduction, and expansion. Outputs an initialized or expanded query vector.
    \item \textbf{Block 2: Document Retrieval:} Handles the parsing, term-frequency vectorization (TF-IDF, N-gram, Local BoW), and initial pruning of the target corpus.
    \item \textbf{Block 3: Ranking:} Maps retrieved documents to queries using either exact lexical dot products or projected semantic spaces (LSA, ESA).
\end{enumerate}
The combination of all parameters resulted in a vast grid-search space. Uniform preprocessing (lowercasing, stopword removal, lemmatization) was enforced across all blocks to guarantee metrics reflected algorithmic choices rather than stemming variations.

% --------------------------------------------------
% 3. Block 1: Query Processing & Expansion
% --------------------------------------------------
\section{Query Refinement and Expansion}
The first block modifies the initial user query representation $q$. We theorized that queries could be optimized either by pruning noisy terms (reduction) or recursively adding semantically relevant neighbors (expansion).

\subsection{Query Reduction}
Query reduction operates under the assumption that non-discriminative terms dilute the cosine matching process. We implemented two specific filters:
\paragraph{IDF Top-$k$ Reduction:} Retains only the $k$ terms with the highest IDF values.
\begin{equation}
score_{idf}(t) = \text{idf}(t)
\end{equation}
With an optimal $k=20$, this yielded negligible gains over the baseline, given that Cranfield queries are inherently concise.

\paragraph{Pseudo-Relevance Feedback (PRF) Pruning:} Explores first-pass retrieval to identify how strongly the query terms are supported by the initial top documents $D_m$:
\begin{equation}
support(t) = \frac{1}{|D_m|} \sum_{d \in D_m} \text{tf}(t, d) \times \text{idf}(t)
\end{equation}
The final score interpolates global rarity with local support:
\begin{equation}
score_{PRF}(t) = \alpha \times \text{idf}(t) + (1 - \alpha) \times support(t)
\end{equation}

\subsection{Matrix-Based Query Expansion}
To counter rigid exact-matching without utilizing the overhead of dense rankers, we formulated an expansion algorithm utilizing pre-computed word-to-word similarity matrices ($S_{w \times w}$) sourced from LSA, ESA, WordNet, and Word2Vec. We apply this explicitly on the \textit{raw term frequencies} rather than IDF matrices to maintain linear scaling behavior during distribution.

To prevent neighborhood noise from overwhelming the original query mass, we implemented a custom dual-stage gating and distribution mechanism:
\begin{enumerate}
    \item \textbf{Dynamic Thresholding:} A base floor is defined by a global similarity quantile. For each term $t_i$, potential expansion neighbors $t_j$ must pass:
    \begin{equation}
    \tau_i = \max\left(\text{floor}, \mu(S_i) + \text{quantile}(S_i) \right)
    \end{equation}
    \item \textbf{L1 Mass Distribution:} Extracted similarities are Min-Max scaled. Crucially, rather than utilizing an exponential Softmax normalization (which skews distribution heavily towards top values), we utilize a proportional L1 normalization to allocate fractional mass from $t_i$ to $t_j$.
    The original term's frequency remains preserved (self-connections are explicitly excluded from distribution loop), meaning actual frequency counts do not diminish; expansion acts strictly additively.
\end{enumerate}

% --------------------------------------------------
% 4. Block 2 & 3: Retrieval and Semantic Ranking
% --------------------------------------------------
\section{Document Retrieval and Latent Ranking}
We test several advanced variations against the standard unigram TF-IDF baseline ($MAP@10 = 0.3061$).

\subsection{Context-Aware Lexical Retrieval (Unordered Local BoW)}
Classical TF-IDF suffers from the phrase-independence assumption. We extended the feature space by generating unordered feature pairs for terms interacting within a sliding local neighborhood window (radius $n$). The scoring augments the baseline document-query matching score $S_{base}$ with context correlation $S_{ctx}$:
\begin{equation}
S(q, d) = S_{base}(q, d) + \beta \cdot S_{ctx}(q, d)
\end{equation}
This reliably outperformed unigram baselines by providing disambiguating context to highly polysemous technical keywords.

\subsection{Latent Semantic Analysis (LSA)}
To solve synonymous wording misalignments, we execute a Truncated Singular Value Decomposition (SVD) on the weighted document TF-IDF matrix $M_{d \times v}$, capturing dimensions of highest variance:
\begin{equation}
M_{d \times v} \approx U_{d \times k} \Sigma_{k \times k} V_{k \times v}^T
\end{equation}
We tuned the latent dimensionality $k=250$. Queries $q_{tfidf}$ are subsequently projected into this space before L2 normalization.
\begin{equation}
q_{latent} = q_{tfidf} \cdot V_k \cdot \Sigma_k^{-1}
\end{equation}
This factorization proved exceptionally strong, ranking overlapping concepts inherently rather than requiring identical string matches.

\subsection{Explicit Semantic Analysis (ESA)}
To target interpretable semantics, ESA utilizes the structural similarity of the test corpus to form an explicit concept space.
It extracts an item-item concept matrix $C$ measuring document self-similarity. A query activates this space forming an overlapping spectrum. ESA also incorporated bigram support $n \in (1,2)$ out of necessity for technical phrases, expanding the feature space to roughly 81,756 dimensions.
\begin{equation}
activations = q_{tfidf} \cdot C^T 
\end{equation}
The 92\% matrix sparsity severely burdened runtime scaling and restricted overall interpretability benefits.

\subsection{Word2Vec Neural Averaging and Soft Cosine}
Dense contextual embeddings natively trained on Cranfield were applied to construct representations.
We found that direct continuous averaging (and even IDF-weighted averaging):
\begin{equation}
v_s = \frac{1}{|s|} \sum_{w \in s} \text{idf}(w) \cdot e_w
\end{equation}
washed out crucial specifics, leading to the worst overall evaluation (MAP 0.2230). To mitigate this, we implemented a Soft Cosine Measure natively combining discrete IDF structure with continuous term distances:
\begin{equation}
\text{SoftCosine}(q, d) = \frac{q^T S d}{\sqrt{(q^T S q)(d^T S d)}}
\end{equation}
Soft Cosine preserved precision relative to flat vectors by anchoring relationships inside the original term matrix.

% --------------------------------------------------
% 5. Experimental Results
% --------------------------------------------------
\section{Experimental Results and Discussion}
The pipeline generates comprehensive IR metrics at the top 10 documents ($k=10$). As seen in Table \ref{tab:results}, standard lexical approaches handle typical exact-matching needs well.

\begin{table}[h]
\centering
\begin{tabular}{lccccc}
\toprule
\textbf{Model Framework} & \textbf{P@10} & \textbf{R@10} & \textbf{MAP@10} & \textbf{nDCG@10} & \textbf{MRR@10} \\
\midrule
TF--IDF (Baseline) & 0.2831 & 0.4094 & 0.3061 & 0.4611 & 0.7461 \\
Query Reduction (PRF-Ext) & 0.2835 & 0.4101 & 0.3067 & 0.4616 & 0.7468 \\
Unordered Local-Context BoW & 0.2925 & 0.4206 & 0.3150 & 0.4721 & 0.7580 \\
Word2Vec (Average Dense) & 0.2298 & 0.3333 & 0.2230 & 0.3649 & 0.6127 \\
Word2Vec (Soft Cosine) & 0.2756 & 0.3990 & 0.2917 & 0.4372 & 0.7018 \\
ESA (Explicit Semantic) & 0.2970 & 0.4430 & 0.3431 & 0.4905 & 0.7910 \\
LSA (Latent Semantic, $k=250$) & 0.3224 & 0.4650 & 0.3703 & 0.5233 & 0.8510 \\
\midrule
\textbf{Hybrid (0.2 TF-IDF + 0.8 LSA)} & \textbf{0.3280} & \textbf{0.4720} & \textbf{0.3795} & \textbf{0.5310} & \textbf{0.8655} \\
\bottomrule
\end{tabular}
\caption{Overall comparative metrics for Cranfield corpus at limit $k=10$.}
\label{tab:results}
\end{table}

The empirical results confirm that LSA effectively supercharges the retrieval pipeline. LSA’s ability to map technically varied descriptions to identical multi-dimensional factors afforded a $22.4\%$ MAP jump. However, pure semantics are often too abstract. Our final proposal pairs exact-matching functionality with dense factors via score weighting:
\begin{equation}
\text{HybridScore}(q, d) = 0.2 \cdot \text{Sim}_{tfidf}(q,d) + 0.8 \cdot \text{Sim}_{LSA}(q,d)
\end{equation}
This hybrid configuration successfully mitigates the penalty semantic networks receive when queries explicitly demand vocabulary strictness. 

% --------------------------------------------------
% 6. Conclusion
% --------------------------------------------------
\section{Conclusion}
The successful integration and automated evaluation of our merging pipeline solidifies the dynamics between lexical rigidity and semantic fluidity in technical domains. By carefully tracing mathematical implementations—from local neighborhood clustering to proportional matrix expansion thresholds—we conclude that on the Cranfield dataset, context-aware lexical algorithms establish a powerful baseline, but Latent Semantic Analysis provides the necessary dimensionality reduction required to conquer vocabulary mismatch. The resulting Hybrid architecture offers a deeply competitive mechanism matching both strict terms and abstract queries evenly.

\end{document}
